Delta-modulated ambient playback refers to a sleep-oriented sound field in which a stable ambient texture is slowly shaped by low-frequency amplitude modulation aligned to the delta range. In practical terms, the device would present a continuous, non-intrusive background soundscape while introducing gentle rhythmic pulses intended to support relaxation and the transition into deep sleep.

The design emphasis is on subtlety rather than stimulation: the carrier audio should remain soft, broadband, and emotionally neutral, while the modulation envelope provides the primary temporal structure. This allows the playback to function as an acoustic scaffold for downshifting arousal without demanding focused attention. In the RBM framework, such a pad is intended to complement endogenous slow-wave dynamics by providing an external rhythmic context that is compatible with sleep onset and maintenance.

A sleep sound pad of this kind would likely be implemented as a bedside speaker or pillow-adjacent transducer with adjustable intensity, modulation depth, and session length. The waveform strategy should prioritize comfort, low spectral harshness, and gradual transitions, so that the user experiences the sound as ambient and continuous rather than as a discrete pulse train. The intended use case is nightly relaxation support, especially for users seeking a low-barrier, non-contact intervention for sleep preparation.
